% ============================================================================
%  chapter12.tex
%  Shared Memory and Memory Mapping
% ============================================================================

\chapter{Shared Memory and Memory Mapping}

\section{The Boldest Abstraction: Memory That Is a File}

Every chapter since the first has treated the address space as a
container of bytes you own. This chapter changes the deal with one
call: \code{mmap}, which \emph{maps} a file---or a region of anonymous
memory---\emph{into your address space}. After the call, reading a
file is just reading memory: \code{p[i]}, not \code{read(fd, ...)}.
And when several processes map the same object as \emph{shared}, they
see the \emph{same physical pages}---shared memory, with no pipe, no
copy, and no kernel involvement in each transfer.

\begin{verbatim}
   The two faces of mmap:

   FILE MAPPING:                    SHARED MEMORY:
   a file's bytes appear in         two (or more) processes map the
   your address space.              same region; they share pages.

        address space               process A         process B
   +----------------------+         +--------+        +--------+
   |  [ p[0] ][ p[1] ]... |         | shared |<======>| shared |
   |    (the file)        |         | pages  |        | pages  |
   +----------------------+         +--------+        +--------+
        ^                             ^ same physical RAM
        | mmap
   +----------------------+
   |   the file on disk   |
   +----------------------+

   Reading the file = dereferencing the pointer.
   Writing the file = assigning through the pointer.
   Sharing between processes = both mapping the same thing SHARED.
\end{verbatim}

To understand why this works---and why it is fast---you need the
virtual-memory model of Chapter~1 in its full form: the address space
is divided into \emph{pages} (4~KB on our machines), and the CPU
translates page-by-page through the \emph{page table}. \code{mmap}
does not copy the file into memory. It merely \emph{points the page
table} at the file's pages. The first access to a page triggers a
\emph{page fault}; the kernel then reads that one page from the file
into RAM, lazily, on demand.

\begin{osconcept}
mmap unifies the two great abstractions of this book---virtual memory
and the file system---into one: \emph{the address space becomes a view
of an object}. The OS maps files lazily, page by page, and shares
physical pages between processes by pointing multiple page tables at
them. Only the kernel can do this, because only the kernel owns the
page tables and the files. Your program just gets a pointer.
\end{osconcept}

\section{mmap: The Call}

\code{mmap} takes six arguments and returns a pointer into your address
space (or \code{MAP\_FAILED}, which is \code{(void *) -1}, not
\code{NULL}). The size must be in bytes, and the kernel works in
pages---a region is rounded up to whole pages whether you like it or
not.

\begin{lstlisting}[caption={Reading a whole file by mapping it.}]
#include <stdio.h>
#include <stdlib.h>
#include <fcntl.h>
#include <sys/mman.h>
#include <sys/stat.h>
#include <unistd.h>

int main(int argc, char *argv[]) {
    int fd;
    struct stat st;
    const unsigned char *p;
    size_t i;

    if (argc != 2) {
        fprintf(stderr, "usage: %s file\n", argv[0]);
        return 1;
    }

    fd = open(argv[1], O_RDONLY);
    if (fd < 0) {
        perror("open");
        return 1;
    }
    if (fstat(fd, &st) < 0) {
        perror("fstat");
        return 1;
    }

    p = mmap(NULL, (size_t) st.st_size, PROT_READ,
             MAP_PRIVATE, fd, 0);
    if (p == MAP_FAILED) {
        perror("mmap");
        return 1;
    }

    /* The whole file is in our address space: just print it. */
    for (i = 0; i < (size_t) st.st_size; i++) {
        putchar(p[i]);
    }

    munmap((void *) p, (size_t) st.st_size);
    close(fd);
    return 0;
}
\end{lstlisting}

The arguments, each doing one job:

\begin{verbatim}
   mmap( addr, length, prot, flags, fd, offset )
        |     |       |     |       |    |
        |     |       |     |       |    +-- start mapping at this file
        |     |       |     |       |         offset (usually 0)
        |     |       |     |       +-- the descriptor, or -1 for
        |     |       |     |            an anonymous mapping
        |     |       |     +-- MAP_SHARED  (visible to others)
        |     |       |          MAP_PRIVATE (private copy, COW)
        |     |       +-- PROT_READ / PROT_WRITE / PROT_EXEC
        |     +-- length in BYTES (rounded up to pages)
        +-- NULL = let the kernel choose the address (almost always)

   munmap(ptr, length)  undoes it, exactly once, when you are done.
\end{verbatim}

\begin{verbatim}
   The lazy mapping, drawn:

      time 0: mmap() succeeds -- NOTHING is loaded yet.
              file:     [p0][p1][p2][p3][p4]
              address:  [    mapped, but empty pages    ]

      time 1: you touch p[0]. The CPU page-faults; the kernel loads
              page 0 from the file into RAM; your read completes.

              file:     [p0][p1][p2][p3][p4]
              address:  [p0][    empty    ]   <- only p0 loaded

      Each first touch of a page is a fault and one disk read.
      Unused pages cost nothing -- the mapping is a promise, not a
      copy. This is why mapping a 1 GB file can succeed instantly.
\end{verbatim}

\begin{osconcept}
The laziness of mmap is demand paging, and it is pure OS machinery.
The page table starts pointing at ``not present,'' the CPU raises a
fault on first touch, and the kernel satisfies that one page from the
file. Your program cannot do this itself; it has no page table. The
OS is betting that you will not touch every byte, and when you do, it
reads only what you actually use.
\end{osconcept}

\section{Anonymous Mappings}

An \emph{anonymous} mapping has no file behind it: the kernel provides
zeroed memory that behaves like an array, and---with
\code{MAP\_SHARED}---can be shared between processes. It is the
modern way to allocate shared memory, and the honest way to get a
big private buffer too (the \code{malloc} of Chapter~13 is built on
it).

\begin{lstlisting}[caption={A shared counter visible to a forked child.}]
#include <stdio.h>
#include <stdlib.h>
#include <sys/mman.h>
#include <unistd.h>
#include <sys/wait.h>

int main(void) {
    long *counter;
    pid_t pid;

    /* One shared page of anonymous memory, writable by all. */
    counter = mmap(NULL, sizeof(long),
                   PROT_READ | PROT_WRITE,
                   MAP_SHARED | MAP_ANONYMOUS, -1, 0);
    if (counter == MAP_FAILED) {
        perror("mmap");
        return 1;
    }
    *counter = 0;

    pid = fork();
    if (pid < 0) {
        perror("fork");
        return 1;
    }

    if (pid == 0) {
        /* Child: add 100 to the SHARED counter. */
        int i;

        for (i = 0; i < 100; i++) {
            *counter = *counter + 1;
        }
        _exit(0);
    }

    /* Parent: wait, then read what the child wrote. */
    waitpid(pid, NULL, 0);
    printf("counter = %ld (child added 100)\n", *counter);

    munmap(counter, sizeof(long));
    return 0;
}
\end{lstlisting}

\begin{verbatim}
   What sharing means here, page by page:

     before fork:   ONE mapping, ONE physical page.
     after fork:    the child inherits the mapping. Because it is
                    MAP_SHARED, BOTH processes' page tables point at
                    the SAME physical page.

     parent: [p] --->  shared page: [ 100 ]  <--- [p] : child
                            ^ same physical RAM

   The child's writes are visible to the parent -- not by copying,
   but because they write the very same bytes. No pipe, no kernel
   buffer, no copy: the fastest possible IPC.
\end{verbatim}

The price of that speed is exactly the theme of this book: shared
writes race. Two processes incrementing the same counter without a
lock lose increments, precisely as two threads did in Chapter~11. The
difference is that now the lock itself must live \emph{in the shared
region}, so both processes can reach it.

\begin{osconcept}
Sharing physical pages is the kernel's greatest privilege: multiple
processes, one piece of RAM, no copying, no serialization. It is also
a deliberate risk, because the kernel cannot police how you use the
bytes---a shared pointer is a shared pointer. The OS provides the
sharing and the synchronization primitives; using them safely is, as
ever, the programmer's contract.
\end{osconcept}

\section{MAP\_SHARED vs MAP\_PRIVATE}

The \code{flags} argument is the one that decides what a \emph{write}
does. This is the crux of the whole mechanism:

\begin{verbatim}
   MAP_PRIVATE   -- the mapping is private to this process.
                     Writes are copy-on-write: the first time this
                     process writes a page, the kernel gives it a
                     PRIVATE COPY, and all further writes touch only
                     the copy. The file (or the other processes) are
                     NOT changed.

   MAP_SHARED    -- the mapping is shared.
                     Writes go straight to the underlying object:
                     the file on disk for a file mapping, or the
                     shared physical pages for everyone else mapping
                     the same region.

   The rule of thumb: PRIVATE to read a file without changing it;
   SHARED to edit a file in place or to communicate between
   processes.
\end{verbatim}

\begin{verbatim}
   COPY-ON-WRITE for MAP_PRIVATE, drawn:

     before any write:                 after the first write to page 0:

     [p0][p1][p2][p3]   (all shared     [p0'](copy)[p1][p2][p3]
        |      with the file)             |        still shared
        +-- the file's pages             +-- this page is now private

   The cost model: reads are free until the pages change; writes cost
   one copy per page, only for pages you actually modify. That is why
   MAP_PRIVATE is the safe default for reading.
\end{verbatim}

For a \emph{shared file} mapping, your writes modify the file's
contents in place---exactly what an editor or a database wants. The
\emph{msync} call pushes the modified pages out to the storage medium
now, rather than waiting for the kernel's background flush. Think of
it as \code{fsync} for a mapping.

\begin{lstlisting}[caption={msync: flush a shared mapping to disk.}]
#include <sys/mman.h>

/* Write the mapped bytes back to the file immediately. */
int flush_mapping(void *addr, size_t length) {
    int r;

    r = msync(addr, length, MS_SYNC);   /* MS_ASYNC = don't wait */
    if (r < 0) {
        return -1;
    }
    return 0;
}
\end{lstlisting}

\begin{osconcept}
The shared/private distinction is kernel policy with teeth. PRIVATE
asks the OS for copy-on-write: the kernel re-points one page at a
fresh copy at the moment of the first write, invisibly and
atomically. SHARED asks the OS to make every write land on the same
physical pages (and, for files, eventually on disk). Choosing between
them is choosing between isolation and communication, and the kernel
enforces that choice per page.
\end{osconcept}

\section{Synchronizing Shared Memory}

Shared memory is fast because there is no kernel in the data path---
which is precisely why there is no kernel in the \emph{coordination}
path either. Two processes that share a counter must synchronize
themselves, with a lock that lives in the shared region. Pthreads
mutexes can do this, but only if told they are \emph{process-shared}:
by default a mutex is only usable by threads of one process.

\begin{lstlisting}[caption={A process-shared mutex guarding shared memory.}]
#include <stdio.h>
#include <stdlib.h>
#include <pthread.h>
#include <sys/mman.h>
#include <unistd.h>
#include <sys/wait.h>

typedef struct shared {
    pthread_mutex_t lock;    /* must live IN the shared region */
    long counter;
} shared;

int main(void) {
    shared *s;
    pthread_mutexattr_t attr;
    pid_t pid;
    int i;

    s = mmap(NULL, sizeof(shared),
             PROT_READ | PROT_WRITE,
             MAP_SHARED | MAP_ANONYMOUS, -1, 0);
    if (s == MAP_FAILED) {
        perror("mmap");
        return 1;
    }

    /* Declare the mutex usable by ANY process (PTHREAD_PROCESS_SHARED). */
    pthread_mutexattr_init(&attr);
    pthread_mutexattr_setpshared(&attr, PTHREAD_PROCESS_SHARED);
    pthread_mutex_init(&s->lock, &attr);
    pthread_mutexattr_destroy(&attr);

    s->counter = 0;

    pid = fork();
    if (pid < 0) {
        perror("fork");
        return 1;
    }

    if (pid == 0) {
        /* Child increments 1000 times, safely. */
        for (i = 0; i < 1000; i++) {
            pthread_mutex_lock(&s->lock);
            s->counter = s->counter + 1;
            pthread_mutex_unlock(&s->lock);
        }
        _exit(0);
    }

    /* Parent increments 1000 times, safely. */
    for (i = 0; i < 1000; i++) {
        pthread_mutex_lock(&s->lock);
        s->counter = s->counter + 1;
        pthread_mutex_unlock(&s->lock);
    }

    waitpid(pid, NULL, 0);
    printf("counter = %ld (expected 2000)\n", s->counter);

    pthread_mutex_destroy(&s->lock);
    munmap(s, sizeof(shared));
    return 0;
}
\end{lstlisting}

\begin{verbatim}
   The layout that makes it work:

     shared region (one anonymous mapping):
     +---------------------------------------------+
     | pthread_mutex_t lock   |  long counter      |
     +---------------------------------------------+
            ^ mapped into BOTH processes' address spaces

     process A:  s->lock, s->counter   (the very same bytes)
     process B:  s->lock, s->counter   (the very same bytes)

   The mutex's memory is shared, so the kernel can see both
   processes contending for it. WITHOUT setpshared(...PROCESS_SHARED),
   the mutex would still be in shared memory but declared private, and
   its behavior would be undefined -- two processes, one lock, no
   contract.
\end{verbatim}

\begin{osconcept}
A mutex is kernel state, and making it process-shared is an explicit
kernel decision (\code{PTHREAD\_PROCESS\_SHARED}) so the OS knows the
lock may be contended across processes, not just threads. The
primitive is the same; only its scope changes. Shared memory without
such a lock is a data race across a process boundary---undefined
behavior with an even bigger blast radius than the thread races of
Chapter~11.
\end{osconcept}

\section{Practical: Editing a File Through a Mapping}

Let us close with the file-backed side of mmap: a program that treats
a file as an array of integers, edits it through the mapping, and
flushes the result to disk. This is the exact technique used by
databases, key-value stores, and editors to modify large files in
place without a single \code{lseek}.

\begin{lstlisting}[caption={mmapfile.c: the file as an array of ints.}]
#include <stdio.h>
#include <stdlib.h>
#include <fcntl.h>
#include <sys/mman.h>
#include <unistd.h>

#define NINTS 8

int main(int argc, char *argv[]) {
    int fd;
    int *vals;
    size_t bytes;
    size_t i;

    if (argc != 2) {
        fprintf(stderr, "usage: %s file\n", argv[0]);
        return 1;
    }

    /* Open (or create) the file read-write. */
    fd = open(argv[1], O_RDWR | O_CREAT, 0644);
    if (fd < 0) {
        perror("open");
        return 1;
    }

    /* Size the file to hold NINTS ints before mapping it. */
    bytes = NINTS * sizeof(int);
    if (ftruncate(fd, (off_t) bytes) < 0) {
        perror("ftruncate");
        return 1;
    }

    /* Map the whole thing, writable and shared. */
    vals = mmap(NULL, bytes, PROT_READ | PROT_WRITE, MAP_SHARED, fd, 0);
    if (vals == MAP_FAILED) {
        perror("mmap");
        return 1;
    }

    /* Write through the mapping: no write(), no lseek. */
    for (i = 0; i < NINTS; i++) {
        vals[i] = (int) (i * i);
    }

    /* Flush now (optional: the kernel would flush eventually). */
    msync(vals, bytes, MS_SYNC);

    printf("wrote %zu ints through the mapping:\n", NINTS);
    for (i = 0; i < NINTS; i++) {
        printf("  vals[%zu] = %d\n", i, vals[i]);
    }

    munmap(vals, bytes);
    close(fd);
    return 0;
}
\end{lstlisting}

\begin{verbatim}
   Why this is a different way to do I/O:

     classic path:   write(fd, buf, n)  -- a syscall, kernel copies
     mmap path:      vals[i] = v;       -- a store instruction; the
                        kernel's page-fault machinery moves the bytes
                        to disk when it must

   The trade-off in one line: mmap removes the copy into and out of
   the kernel for random-access edits, at the price of page-fault
   costs on first touch and giving the kernel control over when the
   data actually reaches disk.
\end{verbatim}

\begin{researchnote}[The theory behind memory-mapped I/O]
mmap is not a trick; it is virtual memory applied to files, and the
theory behind virtual memory is one of computing's foundational
results. In 1970, Peter Denning published ``Virtual Memory''---the
survey that consolidated paging, page tables, and demand loading into
the coherent model this chapter relies on.

\begin{verbatim}
   Denning's framing, condensed:

       "Virtual memory allows programs to run as if they occupied
        far more memory than is physically present, by moving pages
        between main memory and secondary storage on demand."

   Read that with mmap in mind: mapping a 1 GB file with 4 GB of RAM
   is exactly this promise. The file IS the backing store; the page
   fault is the demand; and your pointer is the program's illusion of
   having the whole file in memory.
\end{verbatim}

Denning's survey is a beautifully written summary of why the page
model works, and it explains the behavior of every mmap program you
will ever write---the lazy first touch, the copy-on-write, the
page-sized granularity. When you wonder ``why is my mapping so
fast?,'' the answer was published in 1970.

\paperref{P.~J.~Denning, ``Virtual Memory,'' ACM Computing Surveys
2(3), 1970}{https://dl.acm.org/doi/10.1145/356638.356640}
\end{researchnote}

\section{Structures in Bytes}

The unit of shared memory is the page, and the structure that makes it
work is the page-table entry (PTE)---8 bytes on x86-64.

\begin{verbatim}
   a page: 4096 bytes = 512 eight-byte words.

   a page-table entry, the 8 bytes that name a page:

   bit   0   present           (0 here means "page fault")
   bit   1   writable
   bit   2   user-accessible
   bits  3-11  more flags      (dirty, accessed, ...)
   bits 12-51  physical frame number   <- where the page really lives
   bits 52-62  available to the kernel
   bit  63   NX                (no-execute)

   the four-level walk for a 48-bit address:
   CR3 -> PML4 -> PDPT -> PD -> PT -> 4 KB page
   (each level is a page of 512 eight-byte entries)
\end{verbatim}

\code{mmap} mostly writes entries into these tables: an anonymous
mapping installs ``not present'' entries that fault in zeroed pages; a
\code{MAP\_SHARED} mapping points several processes' PTEs at the same
physical frame number. The ``structure in bytes'' of shared memory is
ultimately this: one 8-byte entry with the same frame number sitting in
two page tables.

\section{Summary}

\begin{verbatim}
   The mapping toolbox, complete:

   mmap(0, len, prot, flags, fd, off)   map a file or anonymous memory
   munmap(ptr, len)                     unmap it exactly once
   MAP_PRIVATE                          copy-on-write: private edits
   MAP_SHARED                           shared pages / edit the file
   MAP_ANONYMOUS                        no file: zeroed shared memory
   MAP_FAILED                           the error sentinel (never NULL)
   page fault                           lazy per-page load on first touch
   msync(ptr, len, MS_SYNC)             flush a shared mapping to disk
   PTHREAD_PROCESS_SHARED               a mutex shared across processes

   Rules that keep mapping code correct:
     1. size a file with ftruncate BEFORE mapping it;
     2. check against MAP_FAILED, not NULL;
     3. pair every mmap with exactly one munmap;
     4. MAP_PRIVATE to read, MAP_SHARED to edit or communicate;
     5. protect shared memory with a process-shared mutex.
\end{verbatim}

\begin{osconcept}
This chapter closed a loop opened in Chapter~1. The address space,
we said, is the process's view of memory; mmap shows it is a view of
\emph{anything the kernel can page}---files, devices, or zeroed
anonymous regions. Sharing physical pages is the fastest and most
dangerous IPC there is, and the kernel provides both the mapping and
the locks, leaving the discipline to you. Nothing in this book
rewards understanding the machine more.
\end{osconcept}

We have built data structures, processes, servers, and shared memory.
The next chapter brings it all into one library: the memory
allocator---the code that hands out and reclaims every byte of the
heap, and the performance cliffs that hide inside \code{malloc}.

